%
% 第10章：训练过程与框架

\clearpage
\cleardoublepage

\chapter{训练过程与框架}

\section{训练循环}

神经网络训练包括在数据上的迭代优化。

\begin{algorithm}[训练循环]
\begin{enumerate}
    \item 初始化模型权重
    \item 对每个epoch：
    \begin{enumerate}
        \item 打乱训练数据
        \item 对每个batch：
        \begin{enumerate}
            \item 前向传播：计算预测
            \item 计算损失
            \item 反向传播：计算梯度
            \item 使用优化器更新权重
        \end{enumerate}
        \item 在验证集上评估
        \item 如果改善则保存检查点
    \end{enumerate}
    \item 返回最佳模型
\end{enumerate}
\end{algorithm}

\section{超参数}

影响训练的关键超参数：

\subsection{学习率}

最重要的超参数。

\begin{itemize}
    \item \textbf{过高：} 发散，损失爆炸
    \item \textbf{过低：} 收敛慢，陷入局部最小值
    \item \textbf{典型范围：} $10^{-5}$ 到 $0.1$
    \item \textbf{Adam默认值：} $10^{-3}$
    \item \textbf{SGD默认值：} $10^{-2}$
\end{itemize}

\subsection{批次大小}

每个训练步的样本数。

\begin{itemize}
    \item \textbf{小批次：} 梯度更噪声，泛化更好
    \item \textbf{大批次：} 梯度更平滑，但需要学习率缩放
    \item \textbf{典型值：} 16, 32, 64, 128, 256
\end{itemize}

\begin{table}[h]
    \centering
    \caption{超参数范围}
    \begin{tabular}{L{2.5cm} L{3cm} L{3.5cm}}
        \toprule
        \textbf{超参数} & \textbf{典型范围} & \textbf{备注} \\
        \midrule
        学习率 & $10^{-5}$ 到 $0.1$ & Adam: $10^{-3}$, SGD: $10^{-2}$ \\
        批次大小 & 16-512 & 取决于GPU内存 \\
        Epochs & 10-1000 & 取决于数据集大小 \\
        权重衰减 & $10^{-6}$ 到 $10^{-2}$ & AdamW: $10^{-2}$典型 \\
        Dropout & 0.1-0.5 & 取决于正则化需求 \\
        动量 & 0.9-0.999 & 用于SGD \\
        \bottomrule
    \end{tabular}
\end{table}

\section{正则化}

防止过拟合的技术：

\subsection{Dropout}

在训练期间随机将激活置零：
\begin{equation}
\mathbf{y} = \mathbf{W} \cdot (\mathbf{x} \odot \mathbf{m})
\end{equation}
其中 $\mathbf{m} \sim \text{Bernoulli}(p)$ 是掩码。

属性：
\begin{itemize}
    \item 子网络集成
    \item 减少共适应
    \item 轻微增加训练时间
    \item 典型dropout率：0.1-0.5
\end{itemize}

\section{反向传播}

通过计算图高效计算梯度的算法。

\begin{algorithm}[反向传播]
\begin{enumerate}
    \item 前向传播：计算所有中间值
    \item 计算损失 $\mathcal{L}$
    \item 反向传播：
    \begin{enumerate}
        \item $\frac{\partial \mathcal{L}}{\partial \mathbf{y}}$ 在输出处
        \item 对每个层（逆序）：
        \begin{enumerate}
            \item 计算 $\frac{\partial \mathcal{L}}{\partial \mathbf{x}} = \frac{\partial \mathcal{L}}{\partial \mathbf{y}} \cdot \frac{\partial \mathbf{y}}{\partial \mathbf{x}}$
            \item 计算 $\frac{\partial \mathcal{L}}{\partial \mathbf{W}}$ 用于权重更新
        \end{enumerate}
    \end{enumerate}
\end{enumerate}
\end{algorithm}

\section{PyTorch框架}

PyTorch是最流行的深度学习研究框架。

\begin{code}
    \captionof{listing}{PyTorch训练循环}
    \label{code:training_chinese}
\begin{minted}{python}
import torch
import torch.nn as nn
import torch.optim as optim
from torch.utils.data import DataLoader

# 模型
model = nn.Sequential(
    nn.Linear(784, 256),
    nn.ReLU(),
    nn.Dropout(0.2),
    nn.Linear(256, 10)
)

# 损失和优化器
criterion = nn.CrossEntropyLoss()
optimizer = optim.Adam(model.parameters(), lr=0.001, weight_decay=0.01)

# 训练循环
device = torch.device('cuda' if torch.cuda.is_available() else 'cpu')
model = model.to(device)

for epoch in range(num_epochs):
    model.train()
    for batch_idx, (data, target) in enumerate(train_loader):
        data, target = data.to(device), target.to(device)
        
        optimizer.zero_grad()
        output = model(data.view(data.size(0), -1))
        loss = criterion(output, target)
        loss.backward()
        optimizer.step()
\end{minted}
\end{code}

\section{CUDA与GPU训练}

PyTorch使用CUDA进行GPU加速。

\begin{code}
    \captionof{listing}{PyTorch中的CUDA操作}
    \label{code:cuda_chinese}
\begin{minted}{python}
import torch

# 检查CUDA可用性
print(torch.cuda.is_available())
print(torch.cuda.device_count())

# 将张量移到GPU
device = torch.device('cuda')
x = torch.randn(1000, 1000).to(device)
y = torch.randn(1000, 1000).to(device)
z = x @ y  # 在GPU上进行矩阵乘法

# 将模型移到GPU
model = MyModel().to(device)

# 混合精度训练
scaler = torch.cuda.amp.GradScaler()
with torch.cuda.amp.autocast():
    output = model(data)
    loss = criterion(output, target)

scaler.scale(loss).backward()
scaler.step(optimizer)
scaler.update()
\end{minted}
\end{code}

\begin{properties}[CUDA原理]
\begin{itemize}
    \item GPU擅长并行操作（矩阵乘法、卷积）
    \item CPU和GPU之间的数据传输成本高
    \item 批处理提高GPU利用率
    \item 混合精度（FP16）减少内存并增加速度
\end{itemize}
\end{properties}

\section{本章小结}

\begin{enumerate}
    \item 训练涉及在批次上的迭代优化
    \item 超参数（学习率、批次大小）关键影响训练
    \item 正则化防止过拟合
    \item 反向传播高效计算梯度
    \item PyTorch提供灵活的深度学习原语
    \item CUDA实现训练加速
\end{enumerate}
